% !TEX encoding = UTF-8
% !TEX TS-program = pdflatex
% !TEX root = ../tesi.tex

%**************************************************************
\chapter{Introduzione}
\label{cap:introduzione}
%**************************************************************
\section{Convenzioni tipografiche}
\label{sec:convenzioni tipografiche}
Durante la stesura del documento sono state adottate le seguenti convenzioni tipografiche:
\begin{itemize}
	\item gli acronimi, le abbreviazioni e i termini ambigui o di uso non comune menzionati vengono definiti nel glossario, situato alla fine del presente documento;
	\item per la prima occorrenza dei termini riportati nel glossario viene utilizzata la seguente nomenclatura: {\color{RoyalBlue}{parola\glsfirstoccur}};
	\item i termini in lingua straniera o facenti parti del gergo tecnico sono evidenziati con il carattere \emph{corsivo}.
\end{itemize}


%**************************************************************
\section{L'azienda}

\myCompany\xspace nasce a Napoli nel 2002 come \textit{software house} ed è rapidamente cresciuta nel
mercato dell'\gls{ICT}, tramutatasi in \textit{System Integrator} e conquistando significative fette di mercato nei settori 
\textit{mobile}, videosorveglianza e sicurezza delle infrastrutture informatiche aziendali. \\
Attualmente, \myCompany ha più di 150 clienti diretti e finali, con un organico aziendale
 di 200 dipendenti distribuiti tra le 5 sedi dislocate in tutta Italia.\\
\myCompany\xspace si pone come obiettivo principale quello di supportare il cliente nella realizzazione,
messa in opera e governance di soluzione IT, sia dal punto di vista tecnologico,
sia nel governo del cambiamento organizzativo.
\begin{figure}[!h] 
    \centering 
    \includegraphics[width=0.5\columnwidth]{logo_azienda.png} 
    \caption{Logo dell'azienda}
\end{figure}

%**************************************************************
\section{Lo stage proposto}
\epigraph{\label{eph: sport}\enquote*{Lo sport dà il meglio di sé quando ci unisce.}}{Frank Deford}
\noindent \productName permette agli utenti di condividere le proprie \gls{will} e partecipare a questa attività sportive.\\
Dal momento che allo stato dell'arte attuale non esiste ancora la suddivisione delle \gls{will}
per gruppi, lo \textit{stage} proposto da \myCompany consiste nell'integrare alla piattaforma già esistente
la visualizzazione delle \gls{will} agli utenti che appartengono agli stessi gruppi.\\  
Gli obbiettivi da raggiungere nel corso dello \textit{stage} sono principalmente due:
\begin{itemize}
    \item sviluppare un \gls{microservizio} utilizzando il \gls{framework} Spring Java per la creazione dei gruppi;
    \item modificare i \glspl{microservizio} esistenti affinché permettano la visualizzazione solo delle \gls{will} di utenti appartenenti allo stesso gruppo.
\end{itemize}
Non è richiesta la modifica del \textit{front end} in modo da adeguarlo alle nuove funzionalità del \textit{back end}, a meno che 
non rimanga tempo da investire su questa attività.

\section{Strumenti utilizzati}
\subsection{Visual Studio Code}
Visual Studio Code è un editor di codice sorgente sviluppato da Microsoft per Windows, Linux e macOS. 
Include il supporto per \textit{debugging}, un controllo per Git integrato, Syntax highlighting, IntelliSense, 
Snippet e \textit{refactoring} del codice.\\
Punto di forza di Visual Studio Code sono le estensioni, grazie alle quali è possibile ampliare 
notevolmente le funzionalità del programma.\\
Le estensioni utilizzate nel corso dello \textit{stage} sono:
    \begin{itemize}
        \item \textbf{GitLens}: amplia le funzionalità di Git integrate in Visual Studio Code;
        \item \textbf{Spring Boot Extension Pack}: raccolta di estensioni per lo sviluppo con Spring Boot;
        \item \textbf{W3C Web Validator}: controlla la validità del \textit{markup} dei documenti HTML e CSS;
        \item \textbf{Web Accessibility}: verifica l'accessibilità dei documenti HTML, evidenziando gli elementi che si potrebbe prendere in considerazione di cambiare e dando suggerimenti su come potrebbero essere modificati;
        \item \textbf{Docker Extension Pack}: raccolta di estensioni per lo gestione dei \textit{container} Docker, Docker images, Dockerfile e file Docker-compose;
        \item \textbf{Angular Extension Pack}: raccolta di estensioni per lo sviluppo con Angular;
        \item \textbf{Extension Pack for Java}: raccolta di estensioni popolari che possono aiutare a scrivere, testare e fare il
        \textit{debugging} di applicazioni Java.
    \end{itemize}
\subsection{Figma}
\label{sub:Figma}
Figma è un \textit{tool} per la progettazione di interfacce, che si rivolge principalmente ai \textit{web designer} che hanno bisogno di un \textit{software} studiato appositamente per realizzare il \textit{design} di siti \textit{web} e applicazioni.\\
Nel contesto dello \textit{stage} è stato utilizzato per la realizzazione delle pagine \textit{web} per la visualizzazione, creazione e modifica di un gruppo. 
\subsection{Git}
\textit{Software} di versionamento utile a tracciare modifiche e cambiamenti di insiemi di file.
\subsection{Postman}
Strumento che permette di eseguire richieste HTTP ad un \textit{server} di \textit{back end}. Quando si lavora con un altro sviluppatore 
\textit{back end} è possibile condividere le \gls{API}, ma la sua vera forza è quella di farci sapere tutto di una richiesta HTTP.
\subsection{DbVisualizer}
DbVisualizer è uno strumento \textit{multi-database} intuitivo e ricco di funzionalità per sviluppatori, analisti e amministratori di \textit{database}, che fornisce un'unica, potente interfaccia su un'ampia gamma di sistemi operativi. \\
DbVisualizer ha un'interfaccia chiara, facile da usare ed è uno strumento che funziona su tutti i principali sistemi operativi, supportando molte varietà di \textit{database}. 


\section{Prodotto ottenuto}
Al termine dello \textit{stage} le integrazioni delle funzionalità con la \textit{web app} sono state realizzata con successo. Tutte le chiamate alle \gls{API} 
sono state testate e sono state integrate anche nel \textit{front end}. 
L'integrazione delle nuove funzionalità permettono alla \textit{web app} di: 
\begin{itemize}
    \item  visualizzare le \gls{will} appartenenti agli utenti che partecipano agli stessi gruppi;
    \item visualizzare le \gls{will} con visibilità globale (quindi che non sono visibili solo agli utenti che partecipano agli stessi gruppi);
    \item visualizzare i gruppi a cui partecipa un utente;
    \item visualizzare e modificare i gruppi creati da un utente;
    \item visualizzare e cercare i gruppi;
    \item creare nuovi gruppi.
\end{itemize}
\section{Organizzazione del testo}

\begin{description}    
    \item[{\hyperref[cap:descrizione-stage]{Il secondo capitolo}}] descrive l'analisi preventiva dei rischi, gli obiettivi dello \textit{stage} e
    la pianificazione delle ore di lavoro.
    \item[{\hyperref[cap:analisi-requisiti]{Il terzo capitolo}}] approfondisce l'analisi dei requisiti del prodotto.
    
    \item[{\hyperref[cap:progettazione-codifica]{Il quarto capitolo}}] approfondisce la fase di progettazione e codifica.
    
    \item[{\hyperref[cap:verifica-validazione]{Il quinto capitolo}}] approfondisce l'accessibilità e la fase di verifica e validazione.
    
    \item[{\hyperref[cap:conclusioni]{Il sesto capitolo}}] contiene un’analisi del lavoro svolto e le conclusioni tratte.
\end{description}
